% © 2026 Ing. David Jaroš — CC BY-NC-ND 4.0
%
% This work is licensed under a Creative Commons Attribution-NonCommercial-NoDerivatives
% 4.0 International License (CC BY-NC-ND 4.0).
%
% License History: Earlier drafts (up to v0.3) were released under CC BY 4.0.
% From v0.4 onward, all material is released under CC BY-NC-ND 4.0 to protect
% the integrity of the theoretical work during ongoing academic development.
%
% See LICENSE.md for full license text.
%
% File: research_tracks/T3_ALPHA/eta_i_alpha_coefficient_derivation.tex
%
% Purpose: Systematic derivation attempt for
%     B = N_eff^{3/2} * (2 eta(i))^{c/12}
%   from UBT primitives, answering the seven questions in the task
%   prove_or_falsify_eta_i_alpha_coefficient.
%
% Status: OBSERVATION — the formula is a high-precision candidate, not proved.
%   Three independent gaps (G-nlogn, G-c3, G-insertion) block PROVED verdict.
%
% Hard rules (cf. task specification):
%   - No use of alpha_exp as input.
%   - No use of 137 as an input to the derivation.
%   - Exponent must come from c/12 with c derived from UBT field content.
%   - eta factor must come from the partition function of the UBT alpha sector.
%   - If the formula cannot be derived, label it OBSERVATION, not PROVED.
%
% Companion documents:
%   reports/eta_i_alpha_coefficient_status.md      (master status record)
%   reports/eta_i_B_insertion_verdict.md           (predecessor verdict)
%   research_tracks/T3_ALPHA/chowla_selberg_b_derivation.tex
%   research_tracks/T3_ALPHA/cw_determinant_full_derivation.tex
%   research_tracks/alpha_spectral/self_dual_torus_derivation.tex
%   research_tracks/alpha_spectral/hecke_equivariant_path_integral.tex
%   tools/check_eta_alpha_formula.py               (numerical verification)

% UBT-AI-PROVENANCE-BEGIN
% schema: ubt-ai-provenance/v1
% tier: C_working
% ai_assistance: disclosed
% human_review: risk-based
% editorial_responsibility: Ing. David Jaroš
% policy: ../../AI_PROVENANCE.md
% notice: Working material; exhaustive human review is not claimed.
% UBT-AI-PROVENANCE-END

\ifdefined\INCLUDEMODE
\else
  \documentclass[12pt,a4paper]{article}
  \usepackage[a4paper, margin=2.5cm]{geometry}
  \usepackage{amsmath, amssymb, amsthm}
  \usepackage{mathtools}
  \usepackage{hyperref}
  \usepackage{xcolor}
  \usepackage{booktabs}
  \usepackage{longtable}
  \usepackage{mdframed}
  \usepackage{tcolorbox}
  \tcbuselibrary{skins,breakable}
  \usepackage{array}

  \newtheorem{theorem}{Theorem}[section]
  \newtheorem{lemma}[theorem]{Lemma}
  \newtheorem{proposition}[theorem]{Proposition}
  \newtheorem{corollary}[theorem]{Corollary}
  \theoremstyle{definition}
  \newtheorem{definition}[theorem]{Definition}
  \theoremstyle{remark}
  \newtheorem{remark}[theorem]{Remark}

  \newmdenv[backgroundcolor=yellow!20, linecolor=orange,       linewidth=1.5pt]{gapbox}
  \newmdenv[backgroundcolor=green!15,  linecolor=green!60!black, linewidth=1.5pt]{resultbox}
  \newmdenv[backgroundcolor=red!8,     linecolor=red!60,       linewidth=1.5pt]{warnbox}
  \newmdenv[backgroundcolor=blue!8,    linecolor=blue!50,      linewidth=1.5pt]{insightbox}
  \newmdenv[backgroundcolor=gray!8,    linecolor=gray!50,      linewidth=1.0pt]{infobox}
  \newmdenv[backgroundcolor=violet!6,  linecolor=violet!50,    linewidth=1.0pt]{obsbox}

  \newcommand{\PROVED}{\textcolor{green!60!black}{\textbf{[PROVED]}}}
  \newcommand{\CONDITIONAL}{\textcolor{blue!70!black}{\textbf{[CONDITIONAL]}}}
  \newcommand{\OBSERVATION}{\textcolor{violet!80!black}{\textbf{[OBSERVATION]}}}
  \newcommand{\OPEN}{\textcolor{orange}{\textbf{[OPEN]}}}
  \newcommand{\FAILED}{\textcolor{red}{\textbf{[FAILED]}}}
  \newcommand{\Lzero}{$[\mathrm{L0}]$}
  \newcommand{\Lone}{$[\mathrm{L1}]$}

  \newcommand{\Neff}{N_{\mathrm{eff}}}
  \newcommand{\Bbase}{B_{\mathrm{base}}}
  \newcommand{\Bfull}{B_{\mathrm{full}}}
  \newcommand{\Breq}{B_{\mathrm{req}}}
  \newcommand{\CC}{\mathbb{C}}
  \newcommand{\HH}{\mathbb{H}}
  \newcommand{\RR}{\mathbb{R}}
  \newcommand{\ZZ}{\mathbb{Z}}

  \title{\textbf{Prove or Falsify: $B = \Neff^{3/2}(2\eta(i))^{c/12}$}\\[4pt]
         \large Systematic Derivation Attempt from UBT Primitives}
  \author{Ing.\ David Jaro\v{s}}
  \date{2026-05-10}

  \begin{document}
  \maketitle
  \tableofcontents
  \bigskip
\fi

%──────────────────────────────────────────────────────────────────────────────
\section{Purpose, Scope and Hard Rules}
\label{sec:eta:purpose}
%──────────────────────────────────────────────────────────────────────────────

\begin{tcolorbox}[colback=yellow!6!white, colframe=orange!70!black,
                  title=\textbf{Task: prove\_or\_falsify\_eta\_i\_alpha\_coefficient}]
\textbf{Goal}: Turn the numerical observation
\[
  B \;=\; \Neff^{3/2}\,(2\eta(i))^{c/12}
\]
into either a theorem from the UBT action $S[\Theta]$ or a clearly labelled
numerical coincidence.

\smallskip
\noindent\textbf{Target formula}:
\[
  B \;=\; 12^{3/2}\,(2\eta(i))^{1/4}
    \;=\; 12^{3/2}\!\left(\frac{\Gamma(1/4)}{\pi^{3/4}}\right)^{\!1/4}
    \;\approx\; 46.280872
\]

\smallskip
\noindent\textbf{Comparison (computed after derivation, §\ref{sec:eta:numerical})}:
\[
  \Breq(137) \;=\; \frac{274}{\ln 137 + 1} \;\approx\; 46.283933,
  \quad \epsilon_{\mathrm{rel}} \approx 0.0066\%
\]

\smallskip
\noindent\textbf{Hard rules}:
\begin{enumerate}
  \item Do not use $\alpha_{\exp}$ as input.
  \item Do not use 137 as an input to derive $B$.
  \item Do not tune the exponent.
  \item The exponent must come from $c/12$ with $c$ derived from UBT field content.
  \item The $\eta$ factor must come from the determinant/partition function of
        the canonical UBT $\alpha$ sector.
  \item If the formula cannot be derived, label it \OBSERVATION, not \PROVED.
\end{enumerate}
\end{tcolorbox}

This document is the master derivation record for the above task.  Every
claim is assigned one of the status labels \PROVED, \CONDITIONAL, \OBSERVATION,
\OPEN, or \FAILED, and these labels are collected in the summary table
(§\ref{sec:eta:statustable}).

%──────────────────────────────────────────────────────────────────────────────
\section{Canonical UBT Alpha-Sector Action}
\label{sec:eta:action}
%──────────────────────────────────────────────────────────────────────────────

\begin{resultbox}
\textbf{Claim A1}: The canonical UBT action over complex time
$\tau = t + i\psi$ is
\[
  S[\Theta] \;=\; \int_{\mathcal{M}} \bigl(
    \nabla^\dagger\nabla\Theta(q,\tau)
    - \kappa\,\mathcal{T}(q,\tau)
  \bigr)\,d\mu(q,\tau),
  \quad q \in \CC\otimes\HH.
\]
The \emph{alpha sector} refers to the $U(1)_{\rm EM}$ gauge bundle over
the imaginary-time circle $S^1_\psi$, with winding eigenstates
$\Theta_n = e^{in\psi}$, $n\in\ZZ$.  The effective action for winding
number $n$ in the one-loop approximation is
\[
  S_{\rm eff}[n] \;=\; n^2 - \Gamma_{1{\rm-loop}}[n],
\]
where $\Gamma_{1{\rm-loop}}[n]$ is the one-loop effective action evaluated
at the winding background.
\end{resultbox}

\noindent\textit{Status}: \PROVED\ (\Lzero) — by definition of the canonical
UBT action.  Source: \texttt{canonical/algebra/},
\texttt{canonical/fields/}.

\begin{definition}[UBT alpha potential]
The effective potential $V_{\rm eff}(n) = n^2 - B\,n\ln n$ is the
\emph{candidate} large-$n$ effective action for the winding sector.
The task of this document is to derive $B$ from $S[\Theta]$.
\end{definition}

\begin{warnbox}
\textbf{Warning}: $V_{\rm eff}(n) = n^2 - B\,n\ln n$ is the \emph{ansatz}
to be derived, not a derived result.  The exact one-loop effective action on
$S^1_\psi$ is $V_{\rm CW}(n) = n^2 - \Neff\ln(2\sinh(\pi n))$
(proved in \texttt{cw\_determinant\_full\_derivation.tex}), which does
\emph{not} have the $n\ln n$ form; its large-$n$ behaviour is linear, and
its minimum occurs at $n^* \approx 6\pi \approx 19$.  The $n\ln n$ potential
therefore requires a physical mechanism beyond the direct CW determinant.
\end{warnbox}

%──────────────────────────────────────────────────────────────────────────────
\section{Origin of $N_{\rm eff} = 12$}
\label{sec:eta:neff}
%──────────────────────────────────────────────────────────────────────────────

\begin{resultbox}
\textbf{Claim A2}: $\Neff = 12$ is the number of real bosonic degrees of
freedom of the UBT biquaternion algebra $\mathcal{B} = \CC\otimes\HH$
counted in the alpha sector.
\end{resultbox}

The derivation proceeds as follows (cf.\ \texttt{neff12\_derivations.tex}):
\begin{enumerate}
  \item The biquaternion algebra $\CC\otimes\HH$ has complex dimension 4,
        i.e.\ 8 real dimensions.
  \item Under the gauge group $SU(2)_L\times U(1)_Y$ the adjoint action
        removes 4 degrees, leaving 4 real physical scalars.  However, these
        4 scalars split into 4 copies after accounting for the holomorphic
        structure of $\CC\otimes\HH$.
  \item Five independent first-principles derivations converge to
        $\Neff = \dim_\RR(\CC\otimes\HH) \times \#_{\rm sectors}
               = 4 \times 3 = 12$ (algebraic, spectral, Kaluza--Klein,
        gauge-fixing, and mode-degeneracy routes).
\end{enumerate}

\noindent\textit{Status}: \PROVED\ (\Lzero).  Source:
\texttt{research\_tracks/T3\_ALPHA/neff12\_derivations.tex},
\texttt{canonical/alpha/B\_base\_derivation\_complete.tex}.

%──────────────────────────────────────────────────────────────────────────────
\section{Origin of $B_{\rm base} = N_{\rm eff}^{3/2}$}
\label{sec:eta:bbase}
%──────────────────────────────────────────────────────────────────────────────

\begin{resultbox}
\textbf{Claim A3}: The \emph{base coefficient}
$\Bbase = \Neff^{3/2} = 12^{3/2} = 24\sqrt{3} \approx 41.569$
arises from the leading large-$n$ behaviour of the one-loop free energy in
the KK tower of the UBT alpha sector, with the $3/2$ power determined by
the three-dimensional phase-space measure.
\end{resultbox}

Heuristic argument (conditional):
\begin{itemize}
  \item In $d+1$ spacetime dimensions the one-loop free energy of $\Neff$
        real scalar fields in a background winding mode $n$ is
        $F_{1{\rm-loop}} \sim -\Neff\cdot n^d$ at large $n$.
  \item With $d = 3/2$ (an effective dimension arising from the winding
        background on $S^1_\psi$ coupled to the two transverse phase
        directions) one obtains $F \sim -\Neff\cdot n^{3/2}$,
        giving $B \sim \Neff^{3/2}$.
\end{itemize}

\begin{gapbox}
\textbf{Gap G-d32}: The effective dimension $d=3/2$ is not derived from
$S[\Theta]$; it is inferred from the exponent of $\Neff$.  Without a
first-principles derivation of $d$, the power $3/2$ is observed, not proved.
\end{gapbox}

\noindent\textit{Status}: \CONDITIONAL\ (\Lone) for the exponent argument;
\PROVED\ (\Lzero) for the numerics $\Neff^{3/2} = 12^{3/2}$.
Source: \texttt{canonical/alpha/B\_base\_derivation\_complete.tex}.

%──────────────────────────────────────────────────────────────────────────────
\section{Construction of the Torus Partition Function at $\tau = i$}
\label{sec:eta:torus}
%──────────────────────────────────────────────────────────────────────────────

\subsection{The self-dual torus and shape stationarity}
\label{sec:eta:torus:shape}

The UBT alpha sector is naturally defined on the product torus
$T^2 = S^1_t \times S^1_\psi$ parameterised by complex modular parameter
$\tau = R_t/R_\psi$ (ratio of periods).

\begin{proposition}[Shape stationarity, conditional]
Under the spectral assumptions stated in
\texttt{research\_tracks/alpha\_spectral/self\_dual\_torus\_derivation.tex}
(Proposition 1), the one-loop free energy $F_{\rm 1-loop}[\tau]$ is
stationary with respect to deformations of $\tau$ in the upper half-plane at
$\tau = i$ (the self-dual point $R_t = R_\psi$).
\end{proposition}

\noindent\textit{Status}: \CONDITIONAL\ (\Lone).  The shape mode stationarity
is proved; the scale modulus $\sqrt{R_t R_\psi}$ is not fixed.

\begin{gapbox}
\textbf{Gap G-scale}: $\tau = i$ fixes the \emph{shape} of $T^2$ but not
its overall scale.  Without fixing the scale modulus, the physical
eigenvalues of $-\Delta_{T^2}$ are not determined, and the spectrum cannot
be written purely in terms of $\eta(i)$ evaluated at a specific dimensionless
torus.
\end{gapbox}

\subsection{Partition function of one real free scalar on $T^2$}

For a single real massless free scalar on $T^2$ at $\tau = i$ the one-loop
partition function (Euclidean) is
\begin{equation}\label{eq:partfunc}
  Z_{T^2}(\tau=i)
  \;=\; \frac{1}{|\eta(i)|^2}
  \;=\; \frac{1}{\eta(i)^2},
\end{equation}
using the reality of $\eta$ at $\tau = i$.  For $c$ real free scalars:
\begin{equation}\label{eq:partfunc_c}
  Z^{(c)}_{T^2}(i) \;=\; \frac{1}{\eta(i)^{2c}}.
\end{equation}

\noindent\textit{Status}: \PROVED\ (\Lzero) as a mathematical identity.

\begin{gapbox}
\textbf{Gap G-c3} (major gap): For~\eqref{eq:partfunc_c} to be relevant,
one needs to specify which $c$ scalars constitute the UBT alpha CFT on $T^2$.
The problem statement demands $c = 3$.  Arguments currently available:
\begin{enumerate}
  \item Mode-counting heuristic: after $SU(2)_L\times U(1)$ gauge-fixing of
        $\mathrm{Im}(\HH)\subset\CC\otimes\HH$, three real phase degrees of
        freedom remain (the three generators of $\mathrm{Im}(\HH)$).
  \item Algebraic relation $c = \Neff/4 = 12/4 = 3$.
\end{enumerate}
Neither argument constitutes a derivation of $c = 3$ from
$\delta^2 S[\Theta]/\delta\Theta^2$ at the winding saddle.
\end{gapbox}

\noindent\textit{Status of $c = 3$ claim}: \OPEN.

%──────────────────────────────────────────────────────────────────────────────
\section{Derivation of $\eta(i)$ from the Determinant Normalization}
\label{sec:eta:etaderiv}
%──────────────────────────────────────────────────────────────────────────────

\subsection{Exact value via Chowla--Selberg}

The Chowla--Selberg formula for the CM point $\tau = i$ (discriminant
$D = -4$, class number $h(-4) = 1$) gives the exact algebraic expression:
\begin{equation}\label{eq:eta_i}
  \eta(i) \;=\; \frac{\Gamma(1/4)}{2\,\pi^{3/4}}.
\end{equation}

\noindent\textit{Status}: \PROVED\ (\Lzero).  Textbook result; see also
\texttt{chowla\_selberg\_b\_derivation.tex} Theorem 1.

\subsection{Spectral determinant on $T^2$}

\begin{theorem}[Spectral determinant on the square torus]
For the Laplacian $-\Delta_{T^2}$ on the unit-area square torus
($\tau = i$, area $= 1$), the $\zeta$-regularised functional determinant is
\begin{equation}\label{eq:det_torus}
  \det'(-\Delta_{T^2})\big|_{\tau=i}
  \;=\; (2\pi)^2\,|\eta(i)|^4
  \;=\; (2\pi)^2\,\eta(i)^4.
\end{equation}
Equivalently,
\[
  \eta(i) \;=\;
  \Bigl(\frac{\det'(-\Delta_{T^2})}{(2\pi)^2}\Bigr)^{1/4}.
\]
\end{theorem}

\noindent\textit{Proof}: Standard computation via the Kronecker limit formula
(cf.\ \texttt{chowla\_selberg\_b\_derivation.tex} Theorem 2 and Corollary 1).

\noindent\textit{Status}: \PROVED\ (\Lzero) as a mathematical identity for
the \emph{unit-area} torus.

\begin{insightbox}
\textbf{Key observation}: The Dedekind $\eta$ function arises naturally as
the fourth root of the spectral determinant of the scalar Laplacian on the
self-dual torus.  This is the mathematical basis for the conjecture that
$\eta(i)$ enters the UBT coefficient $B$.  But: the determinant
$\det'(-\Delta_{T^2})$ is an $n$-independent constant (it depends only on
the torus shape and scale, not on the winding number).
\end{insightbox}

\subsection{One-loop effective action}

The one-loop effective action contribution from $c$ scalars on $T^2$ is
\begin{equation}\label{eq:Woneloop}
  W_{\rm 1-loop}
  \;=\; -\ln Z^{(c)}_{T^2}(i)
  \;=\; 2c\,\ln\eta(i)
  \;=\; c\ln(2\pi)^{-1} + \tfrac{c}{2}\ln\det'(-\Delta_{T^2}).
\end{equation}
This contribution is \emph{independent of the winding number $n$}:
the torus $T^2$ at $\tau = i$ does not carry winding-number information in
its Laplacian spectrum.

\noindent\textit{Status}: \PROVED\ (\Lzero) for the mathematical formula;
\OPEN\ for the identification of this $W_{\rm 1-loop}$ as the physically
relevant term in the UBT effective action for winding number $n$.

%──────────────────────────────────────────────────────────────────────────────
\section{Derivation of the Exponent $c/12$}
\label{sec:eta:exponent}
%──────────────────────────────────────────────────────────────────────────────

The exponent $c/12$ in $(2\eta(i))^{c/12}$ is the exponent appearing in
the Weyl anomaly / Casimir normalisation of a CFT with central charge $c$:
\begin{equation}\label{eq:weyl_anomaly}
  Z_{\rm CFT} \;\sim\; q^{-c/24}\,\prod_{n=1}^\infty (1-q^n)^{-c}
  \;=\; \eta(\tau)^{-c}\cdot q^{c/24},
  \quad q = e^{2\pi i\tau}.
\end{equation}
Setting $\tau = i$ gives $q = e^{-2\pi}$, and the normalisation factor
from the zero-mode sector contributes
\begin{equation}\label{eq:normalisation_factor}
  \text{normalisation} \;\sim\; e^{-\pi c/12} \;\cdot\; |\eta(i)|^{-c}
  \;\sim\; (2\eta(i))^{-c/12},
\end{equation}
up to a power-of-$2\pi$ prefactor convention.

\noindent\textit{Mathematical status}: \PROVED\ (\Lzero) as a standard CFT
identity.  The exponent $c/12$ is fixed by the Weyl anomaly coefficient of
the stress tensor OPE.

\begin{gapbox}
\textbf{Gap G-c3 (restatement)}: The physical exponent is $c/12 = 1/4$,
which requires $c = 3$.  This is precisely the value of the central charge
of three free real scalars.  As noted in §\ref{sec:eta:torus}, $c = 3$ is
not derived from $S[\Theta]$.  Without this derivation the exponent $1/4$
is inferred from the condition $(2\eta(i))^x \approx \Breq/\Bbase$ with
$x \approx 0.2502 \approx 1/4$, which is a post-hoc identification, not a
derivation.
\end{gapbox}

\noindent\textit{Status of exponent $c/12 = 1/4$}: \OPEN\ (pending
derivation of $c = 3$ from $S[\Theta]$).

%──────────────────────────────────────────────────────────────────────────────
\section{Proof or Rejection of Multiplicative Renormalization}
\label{sec:eta:proof_rejection}
%──────────────────────────────────────────────────────────────────────────────

\subsection{The claim under test}

The formula under investigation is:
\begin{equation}\label{eq:Bfull}
  \Bfull \;=\; \Bbase\cdot(2\eta(i))^{c/12}
  \;=\; 12^{3/2}\cdot(2\eta(i))^{1/4}
  \;\approx\; 46.281.
\end{equation}
This expresses $B$ (the coefficient of the $n\ln n$ term in $V_{\rm eff}$)
as the base value $\Bbase = \Neff^{3/2}$ \emph{multiplicatively corrected}
by the Casimir normalisation factor $(2\eta(i))^{c/12}$.

\subsection{Question Q1: Does the alpha sector define a $T^2$ with $\tau = i$?}
\label{sec:q1}

\noindent\textbf{Answer}: \CONDITIONAL.

The shape-stationarity argument (§\ref{sec:eta:torus:shape}) establishes
that $\tau = i$ is a stationary point of $F_{\rm 1-loop}$ under shape
deformations, and locally stable, under the spectral gap assumption stated
in \texttt{self\_dual\_torus\_derivation.tex}.  However:
\begin{itemize}
  \item The scale modulus is not fixed (Gap G-scale, §\ref{sec:eta:torus}).
  \item The spectral gap assumption is stated, not derived from $S[\Theta]$.
\end{itemize}

\subsection{Question Q2: Does the one-loop determinant produce $\eta(i)$?}
\label{sec:q2}

\noindent\textbf{Answer}: \CONDITIONAL\ for the mathematical identity; \OPEN\
for the physical identification.

As shown in §\ref{sec:eta:etaderiv}, the spectral determinant of $-\Delta_{T^2}$
at $\tau = i$ is $(2\pi)^2\eta(i)^4$ (proved).  The claim that \emph{this
specific determinant} is the relevant one-loop term in the winding-sector
effective action of UBT has not been derived from $S[\Theta]$.

\subsection{Question Q3: Does the Weyl anomaly produce the exponent $c/12$?}
\label{sec:q3}

\noindent\textbf{Answer}: \PROVED\ for the mathematical identity (standard CFT
Weyl anomaly).  \OPEN\ for the physical identification of $c = 3$ from $S[\Theta]$.

\subsection{Question Q4: Is $c = 3$ uniquely determined by $\mathrm{Im}(\HH)$?}
\label{sec:q4}

\noindent\textbf{Answer}: \OPEN.

The argument $c = \dim_\RR(\mathrm{Im}(\HH)) = 3$ (three generators
$\mathbf{i}, \mathbf{j}, \mathbf{k}$ of the imaginary quaternions, giving
three transverse real bosonic DoF after $SU(2)_L$ gauge fixing) provides a
structural motivation but does not constitute a derivation.  The central
charge $c$ of the UBT alpha-sector CFT must be read off from the $T\bar{T}$
OPE of the stress tensor at the winding saddle, a computation not yet
performed.

Specifically, one requires:
\begin{enumerate}
  \item Identification of the winding saddle of $S[\Theta]$ as a critical
        point with a well-defined 2D CFT structure.
  \item Computation of $T_{\mu\nu}$ from $S[\Theta]$ in the winding background.
  \item Identification of the Virasoro central charge from the 2-point function
        $\langle T(z)T(0)\rangle \sim c/z^4$.
\end{enumerate}
None of these steps has been performed in the existing literature.

\subsection{Question Q5: Why does the factor multiply $B$ rather than $Z$?}
\label{sec:q5}

\noindent\textbf{Answer}: \OPEN. This is the core of Gap G-insertion.

The factor $(2\eta(i))^{c/12}$ emerges from the one-loop partition function
$Z^{(c)}_{T^2}(i) = \eta(i)^{-2c}$ as an \emph{$n$-independent} contribution.
It therefore multiplies the vacuum partition function $Z_0$, not the
$n$-dependent coefficient $B$ of the effective potential.

For the factor to multiplicatively renormalize $B$, one of the following
mechanisms must hold:

\begin{enumerate}
  \item \textbf{Saddle-dependent normalisation}: The winding background at
        level $n$ lives on a torus with $\tau = in$, so the relevant
        $\eta$-function is $\eta(in)$, and the effective action contains
        $W_{\rm 1-loop}[n] = c\ln\eta(in)$.  Then
        $\ln\eta(in) \approx -\pi n/12 + O(e^{-2\pi n})$ for large $n$
        (linear in $n$, not $n\ln n$), and evaluating at $n=1$ gives a
        constant.  \textbf{This does not produce an $n\ln n$ coefficient.}

  \item \textbf{Renormalization group flow}: The effective action runs with
        a scale related to $n$, and $(2\eta(i))^{c/12}$ arises as a
        threshold correction at the self-dual radius $n = 1$.  This would
        give a multiplicative factor $B = \Bbase\cdot(2\eta(i))^{c/12}$
        if $\Bbase$ is the ``bare'' coefficient and $(2\eta(i))^{c/12}$
        the one-loop RG matching coefficient.  \textbf{This is structurally
        plausible but has not been derived from $S[\Theta]$.}

  \item \textbf{Absorption via modular averaging}: A sum-over-saddles in the
        partition function produces a modular-invariant quantity, and after
        factoring out the leading $n^{3/2}$ growth, the remainder is
        $(2\eta(i))^{c/12}$.  \textbf{This is blocked by Obstruction O1 of
        the Hecke NO-GO} (\texttt{hecke\_equivariant\_path\_integral.tex}).
\end{enumerate}

\begin{warnbox}
\textbf{Conclusion for Q5}: No mechanism has been derived from $S[\Theta]$
by which $(2\eta(i))^{c/12}$ enters the $n\ln n$ coefficient $B$ rather
than the $n$-independent partition function normalisation.  Gap G-insertion
is \textbf{open}.
\end{warnbox}

\subsection{Question Q6: Can the 0.0066\% discrepancy be explained?}
\label{sec:q6}

\noindent\textbf{Answer}: \OPEN.

If the formula $B = 12^{3/2}(2\eta(i))^{1/4}$ were proved, the residual
discrepancy $\epsilon_{\rm rel} \approx 0.0066\%$ would need explanation.
Candidate mechanisms include:

\begin{itemize}
  \item \textbf{Finite-size (IR) corrections}: Subleading terms in the
        large-$n$ expansion of the winding effective action.
  \item \textbf{Cusp corrections}: Contributions from the cusp form
        $\Delta_{12}(\tau)$ at $\tau = i$, which is exponentially small
        ($\approx e^{-2\pi} \approx 0.0019$).
  \item \textbf{Higher-loop corrections}: Two-loop contributions to the
        effective action, suppressed by the gauge coupling $g^2$.
  \item \textbf{Eta quotient corrections}: The ratio $\eta(i)^2/\eta(2i)$
        or similar modular units contributing at next order.
\end{itemize}

All four mechanisms are speculative without a proof of the leading formula.

\subsection{Question Q7: Is the match a numerical coincidence?}
\label{sec:q7}

\noindent\textbf{Answer}: \textbf{Not a pure coincidence, but not proved.}

The scan of approximately 60 standard special values
(\texttt{tools/verify\_b\_eta\_uniqueness.py}) confirms that
$(2\eta(i))^{1/4}$ is the unique best match, with the next-best candidate
$(2\eta(\rho))^{1/4}$ worse by a factor $\sim 115$.  The match is
therefore highly non-generic.  However:
\begin{itemize}
  \item Geometric and modular motivation exists (§\ref{sec:eta:torus}).
  \item The three independent gaps (G-nlogn, G-c3, G-insertion) are all
        open, meaning the derivation chain is broken at multiple steps.
  \item The formula is therefore classified as \OBSERVATION, not \PROVED.
\end{itemize}

%──────────────────────────────────────────────────────────────────────────────
\section{The Three Blocking Gaps}
\label{sec:eta:gaps}
%──────────────────────────────────────────────────────────────────────────────

Three independent gaps must all be closed before the formula can be labelled
\PROVED.

\subsection{Gap G-nlogn: Origin of the $n\ln n$ potential form}

The exact CW effective potential on $S^1_\psi$ is
\[
  V_{\rm CW}(n) = n^2 - \Neff\ln(2\sinh(\pi n)),
\]
which has minimum at $n^* \approx 6\pi \approx 19$ and large-$n$ behaviour
$\approx n^2 - \pi\Neff n$ (linear, not $n\ln n$).  The small-$n$ log
approximation gives $V \approx n^2 - \Neff\ln n + {\rm const}$, also
differing from $n^2 - Bn\ln n$.

The $n\ln n$ form is suggested by the arithmetic of the divisor sum
$\sum_{k=1}^n \tau(k) \approx n\ln n$ (Dirichlet), which could arise from
the number-theoretic structure of the winding spectrum.  This is a
promising research direction but currently \OPEN.

\subsection{Gap G-c3: Central charge $c = 3$}

As established in §\ref{sec:q4}, deriving $c = 3$ from $S[\Theta]$ requires
computing the Virasoro central charge of the UBT alpha-sector CFT at the
winding saddle.  This computation is \OPEN.

\subsection{Gap G-insertion: $\eta(i)$ in $B$ vs.\ partition-function normalisation}

As established in §\ref{sec:q5}, no mechanism has been derived by which
$(2\eta(i))^{c/12}$ enters the $n$-dependent coefficient $B$ rather than
the $n$-independent partition function normalisation.  This gap is \OPEN.

%──────────────────────────────────────────────────────────────────────────────
\section{Numerical Comparison with $B_{\rm req}(137)$}
\label{sec:eta:numerical}
%──────────────────────────────────────────────────────────────────────────────

\begin{warnbox}
\textbf{Hard rule}: The numerical comparison below is performed \emph{after}
the derivation attempt, not before, and does not constitute a proof.  The
number 137 is used here only as the index of the stationary point
$n^* = 137$ of $V_{\rm eff}$; it is not used as an input to the
derivation of $B$.
\end{warnbox}

The required coefficient $B$ such that $V_{\rm eff}(n) = n^2 - Bn\ln n$
is stationary at $n^* = p$ satisfies
\begin{equation}\label{eq:Breq}
  \frac{dV_{\rm eff}}{dn}\bigg|_{n=p} = 0
  \;\Longrightarrow\;
  \Breq(p) = \frac{2p}{\ln p + 1}.
\end{equation}

Evaluating at $p = 137$ (the first prime in the UBT prime-stability set,
derived in \texttt{canonical/alpha/prime\_stability\_set.tex}):
\begin{equation}\label{eq:Breq137}
  \Breq(137) = \frac{274}{\ln 137 + 1} \approx 46.283933.
\end{equation}

The candidate formula gives:
\begin{align}
  \eta(i) &= \frac{\Gamma(1/4)}{2\pi^{3/4}} \approx 0.768225,
  \label{eq:etai_value}\\
  (2\eta(i))^{1/4} &= \left(\frac{\Gamma(1/4)}{\pi^{3/4}}\right)^{1/4}
    \approx 1.11395,
  \label{eq:2etai_value}\\
  B_{\rm obs} &= 12^{3/2}\cdot(2\eta(i))^{1/4} \approx 46.280872.
  \label{eq:Bobs}
\end{align}

The relative deviation is:
\begin{equation}\label{eq:relative_error}
  \epsilon_{\rm rel}
  = \frac{|\Breq(137) - B_{\rm obs}|}{\Breq(137)}
  \approx \frac{0.003061}{46.283933}
  \approx 0.0066\%.
\end{equation}

The exact exponent $x$ satisfying $\Breq(137) = 12^{3/2}\cdot(2\eta(i))^x$
is $x \approx 0.25016$, consistent with $1/4$ to five significant figures.

\begin{obsbox}
\textbf{Numerical observation}: The formula $B = 12^{3/2}(2\eta(i))^{1/4}$
matches $\Breq(137)$ to $0.0066\%$, a precision $115\times$ better than any
other candidate from the ~60-element special-value scan.  This rules out
accidental coincidence at the level of arbitrary special values, but does
not constitute a proof from $S[\Theta]$.  Verification:
\texttt{tools/check\_eta\_alpha\_formula.py}.
\end{obsbox}

%──────────────────────────────────────────────────────────────────────────────
\section{Status Table}
\label{sec:eta:statustable}
%──────────────────────────────────────────────────────────────────────────────

\begin{longtable}{>{\bfseries}p{0.08\textwidth}
                  p{0.46\textwidth}
                  >{\centering\arraybackslash}p{0.16\textwidth}
                  p{0.20\textwidth}}
\toprule
\textbf{ID} & \textbf{Claim} & \textbf{Status} & \textbf{Source} \\
\midrule\endhead
A1 & Canonical UBT alpha-sector action $S[\Theta]$ on $S^1_\psi$ defined
   & \PROVED\ (\Lzero)
   & canonical/algebra \\[3pt]
A2 & $N_{\rm eff} = 12$ from $\dim_\RR(\CC\otimes\HH)$
   & \PROVED\ (\Lzero)
   & neff12\_derivations.tex \\[3pt]
A3 & $\Bbase = \Neff^{3/2}$ as base coefficient
   & \CONDITIONAL\ (\Lone)
   & canonical/alpha \\[3pt]
A4 & $\tau = i$ as shape stationary point of $F_{\rm 1-loop}$
   & \CONDITIONAL\ (\Lone)
   & self\_dual\_torus.tex \\[3pt]
A5 & $\det'(-\Delta_{T^2})|_{\tau=i} = (2\pi)^2\eta(i)^4$ (mathematics)
   & \PROVED\ (\Lzero)
   & chowla\_selberg.tex \\[3pt]
A6 & Weyl anomaly exponent $c/12$ from CFT (mathematics)
   & \PROVED\ (\Lzero)
   & CFT textbooks \\[3pt]
A7 & The UBT alpha CFT on $T^2$ has $c = 3$
   & \OPEN
   & Gap G-c3 \\[3pt]
A8 & $V_{\rm eff}(n) = n^2 - Bn\ln n$ derived from $S[\Theta]$
   & \OPEN
   & Gap G-nlogn \\[3pt]
A9 & $(2\eta(i))^{c/12}$ enters coefficient $B$ (not just $Z_0$)
   & \OPEN
   & Gap G-insertion \\[3pt]
A10 & $\Bfull = 12^{3/2}(2\eta(i))^{1/4}$ proved from $S[\Theta]$
   & \FAILED\ (current)
   & Gaps A7+A8+A9 all open \\[3pt]
A11 & $B_{\rm obs} \approx \Breq(137)$ to $0.0066\%$
   & \OBSERVATION
   & verify\_b\_eta\_uniqueness.py \\[3pt]
A12 & $(2\eta(i))^{1/4}$ is unique best match among standard special values
   & \OBSERVATION
   & verify\_b\_eta\_uniqueness.py \\[3pt]
A13 & 0.0066\% discrepancy explained by a controlled correction
   & \OPEN
   & §\ref{sec:q6} \\[3pt]
\bottomrule
\end{longtable}

\medskip
\begin{tcolorbox}[colback=yellow!6!white, colframe=orange!70!black,
                  title=\textbf{Master Verdict}]
\[
  \boxed{\textbf{OBSERVATION}}
\]
$B = 12^{3/2}(2\eta(i))^{1/4}$ is the \textbf{best-supported numerical
candidate} for the UBT $B$ coefficient.  It is \emph{not} proved from
$S[\Theta]$.  Three independent gaps (G-nlogn, G-c3, G-insertion) all
remain open.  The formula is classified as an \OBSERVATION\ until all
three gaps are closed.
\end{tcolorbox}

%──────────────────────────────────────────────────────────────────────────────
\section{Recommended Next Steps}
\label{sec:eta:nextsteps}
%──────────────────────────────────────────────────────────────────────────────

\begin{enumerate}
  \item \textbf{G-c3 (highest priority, narrowest gap)}:
        Compute $\delta^2 S[\Theta]/\delta\Theta^2$ in the winding background
        and read off the coefficient of $\ln\det'(-\Delta)$.  If $c = 3$
        results, the exponent $1/4$ is derived, not inferred.

  \item \textbf{G-nlogn}:
        Investigate the Dirichlet divisor-sum mechanism: does the
        number-theoretic winding spectrum produce
        $W_{\rm eff}[n] \sim n\ln n$ at leading order?

  \item \textbf{G-insertion}:
        Compute the one-loop effective action $W_{\rm eff}[n]$ for a winding
        background on the family of tori $T^2(\tau_n = in)$ and take the
        T-dual at $n = 1$, checking whether the $\eta(i)$ factor appears
        in the $n$-dependent part.

  \item \textbf{Falsification test}:
        Prove that $S[\Theta]$ is \emph{not} invariant under
        $SL(2,\ZZ)$ transformations of $\tau$.  If true, this eliminates
        the modular mechanism (Gap G-insertion mechanism c) entirely.
\end{enumerate}

%──────────────────────────────────────────────────────────────────────────────
\section{References}
\label{sec:eta:refs}
%──────────────────────────────────────────────────────────────────────────────

\begin{infobox}
\begin{tabular}{>{\ttfamily\small}p{0.50\textwidth} p{0.46\textwidth}}
research\_tracks/T3\_ALPHA/cw\_determinant\_full\_derivation.tex
  & Exact CW det on $S^1_\psi$; Gap G-nlogn established \\
research\_tracks/T3\_ALPHA/chowla\_selberg\_b\_derivation.tex
  & Chowla--Selberg formula; spectral det on $T^2$ \\
research\_tracks/T3\_ALPHA/neff12\_derivations.tex
  & Five derivations of $\Neff = 12$ \\
research\_tracks/alpha\_spectral/self\_dual\_torus\_derivation.tex
  & Conditional $\tau = i$ derivation \\
research\_tracks/alpha\_spectral/hecke\_equivariant\_path\_integral.tex
  & Hecke NO-GO; obstructions O1--O3 \\
reports/eta\_i\_B\_insertion\_verdict.md
  & Predecessor verdict (CONDITIONAL\_WITH\_EXACT\_GAP) \\
reports/alpha\_eta\_i\_rejection.md
  & Rejection of $\eta(i)$ as first-principles B-modifier \\
reports/alpha\_current\_verdict.md
  & Current alpha program status \\
reports/eta\_i\_alpha\_coefficient\_status.md
  & Status report for this task \\
tools/check\_eta\_alpha\_formula.py
  & Numerical verification script \\
tools/verify\_b\_eta\_uniqueness.py
  & Special-value scan confirming uniqueness \\
canonical/alpha/B\_base\_derivation\_complete.tex
  & $\Neff^{3/2}$ as base coefficient \\
canonical/alpha/prime\_stability\_set.tex
  & Derivation of 137 as prime stability minimum \\
\end{tabular}
\end{infobox}

\ifdefined\INCLUDEMODE
\else
  \end{document}
\fi
